\section*{Capítulo \ref{ch:g6:pure-substances-and-mixtures} --- Substâncias puras e misturas}

\begin{solution}{exo:g6:pure-substances-and-mixtures:1}
\omterm{def:g6:pure-substances-and-mixtures:pure-substance}{Substâncias puras}: a água destilada, o bloco de ferro puro, o cristal de
açúcar. \omterm{def:g6:pure-substances-and-mixtures:mixture}{Misturas}: o \omterm{def:g4:air-a-mixture-of-gases:air}{ar}, a água do mar, o leite.
\end{solution}

\begin{solution}{exo:g6:pure-substances-and-mixtures:2}
Homogêneas: a água salgada, a água da torneira. Heterogêneas: o granito,
o óleo sobre a água, a granola.
\end{solution}

\begin{solution}{exo:g6:pure-substances-and-mixtures:3}
Um tipo determinado de substância, com nome e propriedades próprios: por
exemplo a água, o sal, o etanol (e também o açúcar, o oxigênio, o ferro).
\end{solution}

\begin{solution}{exo:g6:pure-substances-and-mixtures:4}
Uma temperatura de ebulição fixa sugere uma \omterm{def:g6:pure-substances-and-mixtures:pure-substance}{substância pura};
\qty{78}{\celsius} é a temperatura de ebulição do etanol.
\end{solution}

\begin{solution}{exo:g6:pure-substances-and-mixtures:5}
Ele contém muitas \omterm{def:g6:pure-substances-and-mixtures:species}{espécies químicas} (água, açúcares, ácidos, vitaminas,
aromas): é uma \omterm{def:g6:pure-substances-and-mixtures:mixture}{mistura}. ``Puro'', na caixa, só quer dizer que nada foi
acrescentado.
\end{solution}

\begin{solution}{exo:g6:pure-substances-and-mixtures:6}
Não: uma \omterm{def:g6:pure-substances-and-mixtures:pure-substance}{substância pura} ferve a uma temperatura fixa. É uma \omterm{def:g6:pure-substances-and-mixtures:mixture}{mistura},
provavelmente água com alguma coisa dissolvida, como a água salgada.
\end{solution}

\begin{solution}{exo:g6:pure-substances-and-mixtures:7}
Um filtro poderia separar a água barrenta (os grãos ficam no papel) e o
suco com polpa (a polpa fica). Ele não consegue separar a água salgada
(o sal está dissolvido) nem o óleo da água (os dois líquidos passam; eles
são separados despejando a camada de cima).
\end{solution}

\begin{solution}{exo:g6:pure-substances-and-mixtures:8}
$78.7 \div 10 = \qty{7.87}{g}$ para \qty{1}{cm^3}: ferro.
\end{solution}

\begin{solution}{exo:g6:pure-substances-and-mixtures:9}
$\frac{357}{478} \approx 0.747$: cerca de \qty{75}{\percent}.
\end{solution}

\begin{solution}{exo:g6:pure-substances-and-mixtures:10}
$2 \times 35 = \qty{70}{g}$ de sais. Porcentagem: $\frac{35}{1000} =
\qty{3.5}{\percent}$.
\end{solution}

\begin{solution}{exo:g6:pure-substances-and-mixtures:11}
Medir as temperaturas de ebulição (cada uma deve ficar fixa durante a
ebulição, e as duas devem ser iguais) e pesar o mesmo volume de cada
líquido (massas iguais). Duas constantes iguais, as duas fixas, indicam
a mesma \omterm{def:g6:pure-substances-and-mixtures:pure-substance}{substância pura}.
\end{solution}

\begin{solution}{exo:g6:pure-substances-and-mixtures:12}
$\frac{10}{90 + 10} = \qty{10}{\percent}$ de etanol. É uma \omterm{def:g6:pure-substances-and-mixtures:mixture}{mistura}: ela
não tem temperatura de ebulição fixa e não aparece na tabela, que só
lista \omterm{def:g6:pure-substances-and-mixtures:pure-substance}{substâncias puras}.
\end{solution}

\begin{solution}{pb:g6:pure-substances-and-mixtures:1}
\textbf{1.} Não. Uma \omterm{def:g6:pure-substances-and-mixtures:pure-substance}{substância pura} e uma \omterm{def:g6:pure-substances-and-mixtures:homogeneous}{mistura homogênea} podem ter
exatamente o mesmo aspecto: é preciso medir alguma coisa.

\textbf{2.} Homogênea: ela é límpida e suas partes não podem ser vistas.

\textbf{3.} Uma constante: uma temperatura de ebulição que fica fixa
enquanto o líquido ferve (ou a massa de um volume conhecido, comparada
com uma tabela).

\textbf{4.} A e B, cujas temperaturas de ebulição ficam fixas, são
\omterm{def:g6:pure-substances-and-mixtures:pure-substance}{substâncias puras}. C é uma \omterm{def:g6:pure-substances-and-mixtures:mixture}{mistura}.

\textbf{5.} A ferve a \qty{100}{\celsius}: água. B ferve a
\qty{78}{\celsius}: etanol.

\textbf{6.} A: $49.8 \div 50.0 = \qty{0.996}{g}$. B: $39.5 \div 50.0 =
\qty{0.790}{g}$. C: $51.2 \div 50.0 = \qty{1.024}{g}$.

\textbf{7.} Sim: cerca de \qty{1.0}{g} para a água, \qty{0.79}{g} para
o etanol.

\textbf{8.} À medida que a água vai embora fervendo, o sal fica para
trás, e o líquido que resta contém cada vez mais sal; ele passa então a
ferver a uma temperatura mais alta.

\textbf{9.} O líquido C contém um sólido dissolvido: é uma solução, uma
\omterm{def:g6:pure-substances-and-mixtures:mixture}{mistura}, como a ebulição dele já tinha mostrado.

\textbf{10.} $100 \times 1.024 = \qty{102.4}{g}$; $\frac{3.5}{102.4}
\approx 0.034$: cerca de \qty{3.4}{\percent} da massa é sal.

\textbf{11.} $\frac{35}{1000} = \qty{3.5}{\percent}$: muito perto da
questão 10.

\textbf{12.} \qty{3.5}{g} em \qty{100}{mL}, logo $3.5 \times 10 =
\qty{35}{g}$ de sal por litro. Pela salinidade e pelo comportamento na
ebulição, C é a água do mar; A é a água destilada e B, o etanol.
\end{solution}
